% !PRINTED 125
\weberpage{125}
%<<CONT
\begin{center}
\S.\,30.\\
\gesperrt{The second Sylow theorem.}
\end{center}

Now let, as above, $P$ be a group of degree $n$ and $p^{\varkappa}$ the highest power of the prime $p$ that divides $n$; let $Q$ be a divisor of $P$ of degree $p^{\varkappa}$, which exists by Theorem I, \S.\,29.

The totality of all elements $c$ of $P$ that satisfy the condition
\weq{1}{c^{-1}Qc = Q,}
among which all the elements of $Q$ are certainly contained, forms a group; for from
\[ Qc = cQ,\quad Qc' = c'Q \]
it follows that $Qcc' = cc'Q$ also holds. These elements $c$ we may call \gesperrt{the elements of $P$ permutable with the group $Q$}. If we denote the group of the $c$ by $R$, then $R$ is a divisor of $P$, and $Q$ a divisor of $R$, and indeed a normal divisor. In particular cases $R$ can be identical with $Q$ as well as with $P$. But if $b$ is any element of $P$ not contained in $R$, then among the elements of the group $b^{-1}Qb$ there is certainly at least one that is not contained in $Q$.

If we denote by $r$ the index of $Q$ with respect to $R$, and by $j$ the index of $R$ with respect to $P$, then
\weq{2}{n = p^{\varkappa}rj,}
and $r$ and $j$ are not divisible by $p$.

The theorem that we have to prove reads:

\gesperrt{II. The index $j$ of $R$ is of the form $pk + 1$, where $k$ is any integer $[j \equiv 1\ (\mathrm{mod}\ p)]$; there are $j$ and no more distinct divisors of $P$ of degree $p^{\varkappa}$, which are all conjugate to one another. Every divisor of $P$ whose degree is a power of $p$ is contained in one of these conjugate groups.}

If $c$ is an element from $R$ of degree $\gamma$, and $s$ is the least positive exponent for which $c^s$ is contained in $Q$, then every other exponent $t$ for which $c^t$ belongs to $Q$ must be divisible by $s$
%CONT>>
